% Приложения. Материал: docs/requirements/traceability.md, docs/research/traceability.md,
% curriculum/lexicon.ru.yaml (транскрипция), docs/research/preregistration.md.
\chapter{Трассировка}\label{app:trace}

\section{Трассировка исследования}

Каждое число главы~\ref{ch:eval} связано с исследовательским вопросом, гипотезой, вопросом GQM, кодом,
который его считает, и рисунком или таблицей, где оно приводится.

\begin{table}[ht]
\centering\small
\caption{Вопрос $\to$ гипотеза $\to$ метрика $\to$ расчёт $\to$ результат.}
\begin{tabularx}{\linewidth}{@{}lllLl@{}}
\toprule
RQ & Гип. & GQM & Метрика (чем считается) & Результат \\
\midrule
RQ-1 & H1.1 & Q1.7 & $\kappa$ «учитель --- учитель» по типу ошибки (чтение / произношение / не ясно) & Т-13 \\
RQ-1 & --- & Q1.8 & $\kappa$ «модель --- учитель» по типу ошибки; доли типов по частям книги & Т-13 \\
RQ-1 & H1.2 & Q1.4 & $\kappa$ «тренажёр --- учитель» против «учитель --- учитель»; $\alpha$ (модуль метрик) & Т-02 \\
RQ-1 & H1.3 & Q1.3 & FAR, FRR, согласие с тренажёром (модуль метрик, SQL-отчёт 05) & Т-01 \\
RQ-1 & H1.4 & Q1.1--Q1.2 & precision, recall, F1 готовых моделей (модуль метрик) & Т-01 \\
RQ-1 & --- & Q1.5 & precision и recall по типу ошибки и букве & Р-08 \\
RQ-2 & с рук. & Q7.1--Q7.3 & лучшая конфигурация agentic loop против базовых подходов: precision, recall, F1, FAR, FRR, парный бутстрэп (оценщик RQ-2); PER на чтении экспертов, если есть дообучение & Т-11, Т-12 \\
RQ-2 & --- & Q7.4--Q7.8 & доля закрытого разрыва, ECE, риск --- покрытие, кривая обучения, вклад источников, решено системой, журнал конфигураций & Т-11, Т-14, Р-14, Р-15 \\
RQ-3 & --- & Q8.1--Q8.4 & инструмент: AUC, лог-потеря, ECE моделей K0--K2; $\rho$ «оценка K2 --- тест учителя, не видящего оценок системы»; подгонка AFM & Т-15, Р-16 \\
RQ-3 & с рук. & Q2.6, Q2.7, Q2.9 & динамика освоения: попыток до усвоения, кривые обучения по навыкам, тест учителя (вход, конец, отложенный), забывание & Р-17, Т-16 \\
RQ-3 & --- & Q2.1--Q2.5, Q2.8 & частоты типов ошибок, путаницы, трудные строки, ошибки по признакам (SQL-отчёты 02--04, 08--11; разведка) & Т-03, Т-04, Р-09, Р-10 \\
E-7 & --- & Q3.1--Q3.2 & изменение балла и повтор ошибки после объяснения (разведка) & Р-11 \\
Метод. & --- & Q6.1--Q6.3 & блокировки по причинам; вмешательства; итерации (журналы) & Р-13, Т-08 \\
\bottomrule
\end{tabularx}
\end{table}

\section{Трассировка требований}

Проект ведёт матрицу, связывающую каждое требование с кодом и тестами; автоматический отчёт заново строит
текущий статус и отпечаток версии кода. На момент написания 59 из 60 прослеженных требований покрыты
проходящими тестами. Полная матрица --- часть репозитория проекта и приводится в электронном приложении.

\chapter{Транскрипция для русскоязычных}\label{app:translit}

Подсказки тренажёра используют упрощённую научную арабско-русскую транскрипцию \cite{SRC-051}, чтобы звуки,
которых нет в русском, оставались различимыми.

\begin{table}[ht]
\centering\small
\caption{Некоторые буквы и их транскрипция в тренажёре.}
\begin{tabular}{@{}cl@{\qquad}cl@{\qquad}cl@{}}
\toprule
Буква & Тр. & Буква & Тр. & Буква & Тр. \\
\midrule
\ar{ح} & х̣ & \ar{ص} & с̣ & \ar{ث} & с̱ \\
\ar{خ} & х & \ar{ض} & д̣ & \ar{ذ} & з̱ \\
\ar{ع} & ‘ & \ar{ط} & т̣ & \ar{ق} & к̣ \\
\ar{ء} & ’ & \ar{ظ} & з̣ & \ar{غ} & г̣ \\
\bottomrule
\end{tabular}
\end{table}

Долгий гласный пишется двумя буквами (аа, ии, уу); слог устроен как «согласная + гласный», \emph{сукун}
закрывает слог.

\chapter{Кратко о предрегистрации}\label{app:prereg}

Предрегистрация (версия 0.5; утверждается и замораживается до сбора данных) фиксирует:
\begin{itemize}
  \item основную гипотезу и три вопроса (RQ-3 --- дополнительный); подтверждающие гипотезы RQ-1
    (таблица~\ref{tab:hyp}), гипотезы RQ-2 и RQ-3 --- с руководителем; обязательные описательные отчёты (журнал
    agentic loop, кривая обучения лучшей конфигурации, вклад источников данных, кривые обучения учеников) и заранее объявленные разведочные
    анализы (E-1, E-4--E-9); всё остальное --- тоже разведочный анализ;
  \item единицу анализа (попытка) и кластеры (ученик);
  \item исключения: попытки заглушки, непригодные записи и «читал не то» (считаются отдельно), попытки до
    заморозки; самоисправления --- отдельным блоком;
  \item участников: взрослые с согласием на запись голоса, обучение модели и тесты учителя; родной язык и опыт
    чтения из анкеты при входе;
  \item разметку: тип ошибки (чтение / произношение / не ясно) обоими учителями на тестовых и каждой десятой
    обучающей попытке; тест учителя, не видящего оценок системы, при входе, в конце пилота и отложенный;
  \item разбиение для RQ-2: около 30\,\% учеников в тесте, список замораживается отпечатком до первого запуска
    agentic loop; разметка теста двумя учителями; подбор конфигураций (кросс-валидация, журнал всех
    конфигураций, фиксированный бюджет) и параметры моделей усвоения --- только на обучающей части; тест
    открывается один раз;
  \item запасной путь: меньше \textasciitilde1\,500 размеченных попыток к 15.04.2027 --- дообучение исключается
    из компонентов agentic loop;
  \item правило остановки: запланированный конец пилота; промежуточные отчёты только описательные;
  \item анализ: интервалы кластерного бутстрэпа, смешанные модели, поправка Холма при $\alpha=0{,}05$;
  \item отложенное: прежние гипотезы H2.1--H4.2 перенесены в план докторантуры (ADR-0064);
  \item заморозку: отпечаток SHA-256 и дата записываются в журнал исследования; любое изменение --- новая
    версия с причиной.
\end{itemize}
